\documentclass[a4paper,12pt]{article}
\usepackage[utf8]{inputenc}
\usepackage[ngerman]{babel}
\usepackage{amsmath}
\usepackage{parskip}

\setlength{\parindent}{0pt}

\title{Modul 293 - Aufgabenblatt 10}
\author{von Lukas Meier}
\date{Unterricht vom 03.03.2026}

\begin{document}

\maketitle

\section*{Aufgabenblatt: HTML-Tabellen (Grundlagen)}

\subsection*{Teil 1: Grundelemente verstehen}
Beantworte die folgenden Fragen schriftlich:
\begin{itemize}
  \item Wofür wird das Element \texttt{table} verwendet?
  \item Was ist der Unterschied zwischen \texttt{th} und \texttt{td}?
  \item Welche Aufgabe hat \texttt{tr}?
\end{itemize}

\subsection*{Teil 2: Tabelle aufbauen}
Erstelle eine Tabelle mit:
\begin{itemize}
  \item einer Kopfzeile mit den Spalten \texttt{Name}, \texttt{Fach}, \texttt{Note}
  \item mindestens vier Datenzeilen
  \item sinnvoller Reihenfolge und korrekter Einrückung
\end{itemize}

\subsection*{Teil 3: Struktur verbessern}
Erweitere deine Tabelle aus Teil 2 mit:
\begin{itemize}
  \item \texttt{caption}
  \item \texttt{thead}
  \item \texttt{tbody}
\end{itemize}

Notiere kurz, warum diese Elemente nützlich sind.

\subsection*{Teil 4: Zellen zusammenfassen}
Erstelle eine zweite Tabelle (z.\,B. Stundenplan), in der du verwendest:
\begin{itemize}
  \item mindestens einmal \texttt{colspan}
  \item optional einmal \texttt{rowspan}
\end{itemize}

Beschreibe in 2--3 Sätzen, was durch das Zusammenfassen passiert.

\subsection*{Teil 5: Fehler finden und korrigieren}
Der folgende Code enthält Fehler:

\begin{verbatim}
<table>
  <thead>
    <th>Produkt</th>
    <th>Preis</th>
  </thead>
  <tbody>
    <tr>
      <td>Tastatur
      <td>49.-</td>
    </tr>
  </tbody>
</table>
\end{verbatim}

\begin{itemize}
  \item Finde mindestens drei Fehler.
  \item Korrigiere den Code vollständig.
  \item Begründe jede Korrektur in einem Satz.
\end{itemize}

\subsection*{Teil 6: Praxisaufgabe}
Erstelle eine kleine HTML-Seite mit einem Bereich \texttt{main}, der eine sauber strukturierte Tabelle enthält.

Anforderungen:
\begin{itemize}
  \item Tabelle mit \texttt{caption}, \texttt{thead}, \texttt{tbody}
  \item mindestens 5 Datenzeilen
  \item mindestens eine zusammengefasste Zelle mit \texttt{colspan}
  \item lesbarer, eingerückter Code
\end{itemize}

\subsection*{Teil 7: Reflexion}
Beantworte die folgenden Fragen kurz:
\begin{itemize}
  \item Was war heute bei Tabellen am einfachsten?
  \item Wo gab es Unsicherheiten?
  \item Woran erkennst du, dass eine HTML-Tabelle korrekt strukturiert ist?
\end{itemize}

\end{document}
